\section{Теоретична основа}
\label{sec:theory}

Предходната глава позиционира работата спрямо съществуващите защити. Настоящата
глава въвежда само понятията, необходими за разбирането на собствения принос:
формално определение на инжекцията, трите принципно различни начина за нейното
предотвратяване, разликите между езиците за заявки на изследваните хранилища и
ролята на времето на компилация като граница на доверие. Последната точка е
мостът между определението и реализацията, разгледана в
глава~\ref{sec:enforcement}.

\subsection{Инжекцията като промяна на синтактичната структура}
\label{sec:theory-definition}

Интуитивното описание на инжекцията („въведеният текст се изпълнява като код“) е
недостатъчно точно, за да служи за критерий при оценка. В настоящата работа се
използва определението на Су и Васерман~\cite{su2006essence}. Заявката се
разглежда като изречение от формален език и съответно като синтактично дърво.
Всеки фрагмент от заявката, произлизащ от недоверен вход, трябва да попада изцяло
в едно листо на това дърво, тоест да бъде стойност. \textbf{Инжекция} е налице
тогава, когато недоверен вход променя синтактичната структура на заявката: когато
той обхваща повече от едно листо или въвежда нови вътрешни възли.

Определението има две предимства. Първо, то е независимо от конкретния полезен
товар (\textit{payload}): то не изисква списък от опасни символи, а критерий за
структурна цялост. Второ, то е независимо от езика. Едно и също понятие обхваща
SQL, документните заявки на MongoDB, Cypher, CQL и командния протокол на Redis,
стига за всеки от тях да е известна граматиката, по която хранилището разбира
заявката. Именно тази независимост позволява изследователският въпрос да бъде
поставен за цяло семейство от класове инжекции, а не за всеки поотделно.

Класификацията на Халфонд, Виегас и Орсо~\cite{halfond2006classification}
допълва определението от страна на практиката. Тя разграничава атаките по
механизъм (тавтологии, обединения чрез \code{UNION}, натрупани заявки, заявки,
предизвикващи грешка, слепи атаки по истинност и по време) и по намерение
(извличане на данни, заобикаляне на удостоверяване, промяна на данни, отказ от
услуга). Корпусът от полезни товари, използван в
глава~\ref{sec:evaluation}, е подреден по същите механизми.

\subsection{Три начина за отделяне на данните от кода}
\label{sec:theory-remedies}

Общата причина за всички инжекции е една: данни и код се пренасят по един и същ
канал, най-често като един низ, и хранилището ги разделя едва при синтактичния
анализ. Всяка защита следователно отговаря на въпроса как недоверената стойност
да остане листо в дървото. Принципно се различават три подхода.

\textbf{Екраниране} (\textit{escaping}). Стойността остава в същия канал, но се
преобразува така, че синтактичният анализатор да не я разпознае като граница на
листото. Сигурността зависи от пълното и правилно преобразуване за конкретния
диалект и кодиране на символите. Пропуск на един специален символ или разлика
между представата на приложението и на хранилището за кодирането е достатъчна,
за да се загуби гаранцията.

\textbf{Свързване на параметри} (\textit{parameter binding}). Заявката се изпраща
с маркери на местата на стойностите, а самите стойности пътуват отделно, по
протокол, който хранилището никога не подлага на синтактичен анализ като заявка.
Стойността не може да стане вътрешен възел, защото не участва в анализа изобщо.
Това е механизмът на подготвените заявки~\cite{owasp_sqli_prevention}.

\textbf{Конструиране} (\textit{construction}). Структурата на заявката не се
сглобява от низове, а се изразява в средство, в което недоверените данни нямат
път към структурата. Свързването защитава стойностите, но оставя отворен въпроса
кой определя структурата. Конструирането отговаря на този въпрос.

Разграничението е съществено, защото свързването не покрива всички части на
заявката. Нито едно от изследваните хранилища не позволява параметър на мястото
на име на таблица или колона. Затова приложения, които използват подготвени
заявки за стойностите, нерядко сглобяват идентификаторите като низове, и именно
там остава уязвимостта. Настоящата работа изследва конструирането като допълнение
към свързването, което затваря тази празнина.

\subsection{Езиците на изследваните хранилища}
\label{sec:theory-languages}

Определението от точка~\ref{sec:theory-definition} предполага граматика. Петте
класа инжекции в обхвата на работата се различават именно по това каква
граматика анализира хранилището и дали изобщо анализира текст, получен от
приложението.

\begin{table}[H]
\centering
\caption{Граматика и клас инжекция при изследваните хранилища}
\label{tab:theory-languages}
\begin{tabularx}{\textwidth}{L{2.6cm}L{3.2cm}Y}
\toprule
\textbf{Хранилище} & \textbf{Какво се анализира} & \textbf{Клас инжекция} \\
\midrule
PostgreSQL, MySQL, SQLite & текст на SQL & SQL инжекция: тавтологии, \code{UNION}, натрупани заявки, слепи атаки \\
Cassandra & текст на CQL & CQL инжекция, сходна със SQL, но с по-беден език \\
Neo4j & текст на Cypher & Cypher инжекция, включително извикване на процедури \\
MongoDB & документ BSON, не текст & инжекция на оператори: стойност, която се интерпретира като {\ttfamily \$ne} или {\ttfamily \$where} \\
Redis & кадри RESP с префикс за дължина & инжекция на протоколно ниво чрез \code{CR LF}, ако командата се сглобява като текст \\
\bottomrule
\end{tabularx}
\end{table}

Таблица~\ref{tab:theory-languages} показва, че хранилищата не са еднородни по
отношение на заплахата. При SQL, CQL и Cypher хранилището анализира текст и
опасността е класическата. При MongoDB заявката е структуриран документ и
инжекцията се изразява по друг начин: стойност, която приложението очаква да е
низ, пристига като обект с ключ, започващ с {\ttfamily \$}, и се интерпретира като
оператор. При Redis протоколът RESP предава всеки аргумент с изрично указана
дължина. Ако клиентът кодира аргументите правилно, съдържанието на аргумента не
може да бъде прочетено като граница на команда, независимо от символите в него.
Следователно за две от седемте хранилища правилното кодиране премахва цялата
граматика, от която атакуващият би могъл да излезе, а не само я екранира.

\subsection{Времето на компилация като граница на доверие}
\label{sec:theory-compile-time}

Конструирането от точка~\ref{sec:theory-remedies} изисква средство, в което
недоверените данни нямат път към структурата. В езика C++ такова средство е
разграничението между стойностите, известни по време на компилация, и тези,
известни едва по време на изпълнение.

Аргументите на шаблон, които не са типове (\textit{non-type template
parameters}), трябва да бъдат \textbf{константни изрази}. Стойност, която
зависи от входа на програмата, по дефиниция не е константен израз и не може да
бъде такъв аргумент. Ако следователно всеки елемент на синтактичната структура
на заявката (име на таблица, име на колона, вид на операцията, оператор на
сравнение, наличие на клауза) е закодиран като аргумент на шаблон или като тип,
то никоя стойност, получена по време на изпълнение, не може да го промени. Входът
може да достигне заявката само там, където типовата система изрично го допуска.

Аргументът може да се формулира като свързване на определението с типовата
система:

\begin{enumerate}
    \item Нека всички вътрешни възли на синтактичното дърво на заявката са
          определени от аргументи на шаблони и типове. Тогава те са фиксирани в
          момента на компилация.
    \item Нека единственият път, по който стойност по време на изпълнение достига
          заявката, е типизиран заместител, който се свързва като параметър
          съгласно точка~\ref{sec:theory-remedies}.
    \item Тогава недоверен вход може да заеме само листо и не може да промени
          структурата. Съгласно определението на Су и
          Васерман~\cite{su2006essence} инжекция не е възможна.
\end{enumerate}

Двете предпоставки не се приемат на доверие. Първата се проверява в
глава~\ref{sec:evaluation} чрез програми, които се опитват да подадат стойност по
време на изпълнение на мястото на идентификатор или структура и трябва да бъдат
отхвърлени от компилатора. Втората се проверява чрез диференциален тест и чрез
анализа на всеки конектор в глава~\ref{sec:enforcement}. Именно проверката на
втората предпоставка разкри случай, в който тя не е била изпълнена: казусът в
глава~\ref{sec:case-study}.

Цената на подхода е симетрична на гаранцията. Ако идентификаторите и структурата
не могат да зависят от входа, то и програмистът не може да ги направи зависими.
Схеми, които се определят динамично, не се поддържат. Това не е пропуск, а
същността на гаранцията, и е обсъдено като ограничение в
глави~\ref{sec:threat-model} и в заключението.
